\section{From 929 to 6556: The Scroll Construction}

This is the most constructed object in the paper, and for that reason I want its steps to be unusually transparent. It is not a literal divine name. It begins with the conventional chapter total 929, converts that value to base 12, and then applies a one-digit palindromic closure.

\subsection{Construction from the Chapter Total}

Under the chapter-counting convention now common in printed Tanakh editions, the total is 929. The 929 Project notes that this division arose in the late Middle Ages, is conventionally associated with Stephen Langton, and was later adopted in Jewish use.\cite{project929,verse_division} The number is therefore exact within a received chapter system, not an ancient division internal to the biblical text. The first mathematical step is:

\[
929_{10}=655_{12},
\]

because

\[
6(12^2)+5(12)+5=864+60+5=929.
\]

The decimal source value 929 is itself palindromic. Its base-12 form, 655, has the pattern $abb$: the last two digits are already equal. A one-digit right extension has the form $655x$. For it to read the same backward, the new last digit must equal the first digit. Thus $x=6$, and no other appended digit works:

\[
655_{12}\xrightarrow{\text{one-digit palindromic closure}}6556_{12}.
\]

The two arrows in the full sequence do different work:

\[
929_{10}\xrightarrow{\text{base conversion}}655_{12}
\xrightarrow{\text{palindromic extension}}6556_{12}.
\]

The second arrow is not another conversion. It is a defined closure operation in which the required digit 6 converts the $6\,5\,5$ pattern into the reversal-invariant string $6\,5\,5\,6$. The completed string contains the original 655 in its first three places and its reversal 556 in its last three, overlapping on the central 55. That forward/reverse overlap is exact, although it follows automatically from the palindromic closure rather than supplying a separate test.

\subsection{What the Canonical Label Means}

Earlier drafts called the result a ``22-scroll constant.'' The historical basis for 22 is Josephus, who speaks of twenty-two sacred \emph{books} in \emph{Against Apion} 1.38--41.\cite{josephus22} He does not provide the detailed scroll arrangement implied by the later shorthand, while the standard modern Jewish enumeration is 24 books.\cite{bible_canon24} I therefore use \emph{scroll construction} for the mathematical object and treat its contact with 22 as a result of the collapse path, not as proof of one uniquely fixed ancient scroll count.

\subsection{Exact Value and Symmetry}

The decimal value of the completed base-12 string is

\begin{align*}
6556_{12}
&=6(12^3)+5(12^2)+5(12)+6\\
&=10368+720+60+6\\
&=11154_{10}.
\end{align*}

Its palindrome status is exact:

\[
\operatorname{reverse}(6556)=6556.
\]

This is the strongest strict property of the completed construction. It links a decimal palindrome at the source, $929_{10}$, to a base-12 palindrome after conversion and one-digit closure.

\subsection{The Frame and the Finite Check}

The completed notation has an exact outer/inner partition:

\[
(6+6)+(5+5)=12+10=22.
\]

The outer digits therefore sum to 12. The inner pair is also a complete admitted value rather than merely two matching digits:

\[
\text{Adonai}=65_{10}=55_{12}.
\]

This gives the construction an Adonai center while the whole notation passes through 22 and then 4. These statements describe the displayed structure; they do not mean that Adonai generates the chapter count.

The pattern was checked exhaustively against every three-digit base-12 integer from $100_{12}$ through $\mathrm{BBB}_{12}$. There are 1,584 such integers. Exactly 132 have the form $abb$ and therefore admit a one-digit right closure $abba$. Among those 132 closures, 12 have outer digits summing to 12, 11 have all four digits summing to 22, and only one has both properties: $6556_{12}$. Five converted decimal values have a first digit sum of 12; only $6556_{12}=11154_{10}$ combines that decimal-value sum with the notation sum of 22.

This is a complete finite count, but it is not a prospective probability. The outer/inner partition and the comparison were articulated after $6556_{12}$ was already known. The scan establishes that the conjunction is unique within the declared closure family; it does not establish why that conjunction occurs.

\subsection{Two Distinct Collapse Paths}

Earlier wording assigned the construction directly to the 12-collapse cluster. That description combined two different operations.

\paragraph{Base-12 notation-collapse.}
Summing the displayed base-12 digit values gives

\[
6+5+5+6=22_{10},
\]

followed by

\[
2+2=4.
\]

Thus the notation-level path is

\[
6556_{12}\longrightarrow22_{10}\longrightarrow4.
\]

It passes through decimal 22 and terminates at 4 under the manuscript's collapse rule. It does not terminate at 12.

\paragraph{Decimal-value collapse.}
If the base-12 number is first converted to its decimal value, a separate path appears:

\[
6556_{12}=11154_{10}
\longrightarrow1+1+1+5+4=12.
\]

The completed construction therefore has a legitimate 12 endpoint, but only through collapse of the converted decimal value. The manuscript must name this basis whenever the 12 result is used.

\subsection{The Intermediate 22 and the Final 4}

The notation-collapse's intermediate value, $22_{10}$, makes direct contact with the 22-letter alphabet count used in the model. It also reproduces the visible digit string in YHWH's base-12 form:

\[
\text{YHWH}=26_{10}=22_{12}.
\]

These two 22s are not numerically equal:

\[
22_{10}=1\mathrm{A}_{12},
\qquad
22_{12}=26_{10}.
\]

The alphabetic link is therefore exact at decimal 22, while the YHWH link is a cross-basis digit-string echo. Keeping those claims distinct makes the observation more precise, not less meaningful.

The final collapse to 4 is also structurally relevant. At the interpretive level, it can be read against the four letters of the Tetragrammaton. The calculation itself establishes only the endpoint 4. The Tetragrammaton reading is a theological interpretation of that endpoint, but it is a natural one within the manuscript's identity framework and should be retained as such.

The complete notation path can therefore be summarized as a three-stage structure:

\begin{center}
palindromic identity $\longrightarrow$ alphabetic 22 $\longrightarrow$ fourfold endpoint.
\end{center}

\subsection{Relations to the Divine-Name Corpus}

The construction has three especially close contacts with audited rows.

\paragraph{Adonai.}
The central pair is $55_{12}$, the complete base-12 value of Adonai. It is also the overlap shared by the forward source string 655 and the reversed suffix 556. This is an exact same-basis containment inside the constructed palindrome, not an equality between Adonai and the complete object.

\medskip

\paragraph{YHWH.}
Both $22_{12}$ and $6556_{12}$ are base-12 palindromes. The former is an exact conversion of YHWH's standard gematria; the latter is a canonical-scale construction. This supports a micro/macro analogy, provided it is described as an analogy between palindromic signatures rather than an arithmetic derivation of one from the other.

\medskip
\paragraph{HaKadosh Baruch Hu.}
The full article form computes to a standard value of $655_{10}$, while the chapter total gives $655_{12}$. These are different numbers that share the same displayed digit string. The chapter-derived $655_{12}$ then serves as the prefix of $6556_{12}$. This is a potentially important cross-basis echo, but the manuscript should not write $655_{10}=655_{12}$ or imply that the title generates the chapter total.

\subsection{Evidential Status}

The argument is clearest when its levels are summarized explicitly:

\begin{description}[leftmargin=4.3cm,style=multiline]
\item[Exact arithmetic] $929_{10}=655_{12}$ and $6556_{12}=11154_{10}$.
\item[Exact symmetry] $929_{10}$ and $6556_{12}$ are palindromes in their stated bases.
\item[Defined construction] Appending 6 is the one-digit palindromic closure of $655_{12}$.
\item[Finite comparison] Among 132 eligible closures, $6556_{12}$ alone combines a notation sum of 22 with a converted-value sum of 12; it also alone combines outer sum 12 with notation sum 22.
\item[Exact collapse] The notation path is $22_{10}\rightarrow4$; the decimal-value path is $11154\rightarrow12$.
\item[Interpretation] The inner $55_{12}$ is Adonai, the 22 evokes the alphabet, the visible 22 string echoes YHWH, and the 4 evokes the four-letter Name.
\item[Historical basis] The 929 count uses the received medieval chapter division; the 22-unit reference is attested as a book count in Josephus.
\item[Intentionality question] The arithmetic does not by itself establish deliberate encoding.
\end{description}

\subsection{Corrected Formulation}

The scroll construction is a base-12 palindrome produced by a transparent closure operation on the base-12 form of 929. The appended 6 is forced by the one-digit palindrome requirement. Its outer digits sum to 12, its inner digits form Adonai's $55_{12}$, its notation-collapse passes through decimal 22 and ends at 4, and its converted decimal value collapses separately to 12. The finite scan makes the conjunction unique within the declared 132-member closure family, while the base labels, post-discovery status, and historical conventions remain visible.

What remains is the observation I wanted to preserve: a compact constructed object that joins chapter-count conversion, a forced palindromic closure, the complete value of Adonai at its center, an outer sum of 12, an alphabetic intermediate, a fourfold endpoint, and the separate exact decimal sum $1+1+1+5+4=12$. The construction should not be mistaken for something discovered ready-made in the chapter total, but neither should its final collapse to 4 be discarded as meaningless. The arithmetic is clear enough for the reader to decide how much weight the structure can bear.
